% Revize 2026-10-28: obsah, matematika a návaznost; podrobnosti a časy
% viz planovani/REVIZE_PRIPRAV_1_ROCNIK.md.
% Aktuální porovnání se skutečným obsahem místních PDF Realisticky.cz:
% externi_zdroje/realisticky/01_Zakladni_poznatky/06_Mocniny/010601_Mocniny_s_prirozenym_mocnitelem/010601_Lekce.pdf
% Revize 2026-10-27: cílené doplnění; nové úlohy a řešení jsou vlastní tvorba.
% Metodický vzor: 4_rocnik/pripravy/Uvod_do_kombinatoriky.pdf (samostatný dokument).
% Hodina 25, 1. ročník SPŠ. Rozsah: Rozpis_vyuky_MAT_1_rocnik.tex, řádek 25.

% Plán: 1_rocnik/tematicke_plany/SPS.pdf; společná projektová šablona.

% Učebnice: ucebnice/01_Zakladni_poznatky/Mocniny_a_odmocniny/Mocniny a Odmocniny.tex

% Původ úloh: P1: učebnice, výklad znamének; P2 a procvičení: „Rychlé mocniny – výpočet z paměti“ a „Pravdivost tvrzení“, vybrané položky; poslední položka vlastní.

% Metodický zdroj: externi_zdroje/realisticky/01_Zakladni_poznatky/06_Mocniny/010601_Mocniny_s_prirozenym_mocnitelem/010601_Lekce.pdf

% URL: http://www.realisticky.cz/ucebnice/01%20Matematika%20S%C5%A0/01%20Z%C3%A1kladn%C3%AD%20poznatky/06%20Mocniny/01%20Mocniny%20s%20p%C5%99irozen%C3%BDm%20mocnitelem.pdf

% Metodický zdroj: externi_zdroje/realisticky/01_Zakladni_poznatky/06_Mocniny/010602_Vzorce_pro_mocniny_I/010602_Lekce.pdf

% URL: http://www.realisticky.cz/ucebnice/01%20Matematika%20S%C5%A0/01%20Z%C3%A1kladn%C3%AD%20poznatky/06%20Mocniny/02%20Vzorce%20pro%20mocniny%20I.pdf

\documentclass[11pt,a4paper]{article}

% Spolecny styl pro ucitelske pripravy MAT (Work/MAT)
% Verze 1.1, 2026-10-23
\NeedsTeXFormat{LaTeX2e}
\ProvidesPackage{priprava}[2026/10/23 v1.1 Sablona ucitelskych priprav MAT]

% Zakladni typografie a cestina
\RequirePackage[T1]{fontenc}
\RequirePackage[utf8]{inputenc}
\RequirePackage[czech]{babel}
\RequirePackage{lmodern}
\RequirePackage{microtype}
\RequirePackage{amsmath,amssymb}
\RequirePackage{xcolor}
\RequirePackage{enumitem}
\RequirePackage[most]{tcolorbox}
\RequirePackage{tabularx,booktabs,array}
\RequirePackage{fancyhdr}
\RequirePackage{needspace}
\RequirePackage{geometry}
\RequirePackage[hidelinks]{hyperref}

\geometry{a4paper,top=17mm,bottom=18mm,left=23mm,right=23mm}
\definecolor{notegray}{gray}{0.39}
\setlength{\parindent}{0pt}
\setlength{\parskip}{2.5pt plus 1pt minus 0.5pt}
\setlist[itemize]{topsep=2pt,itemsep=1pt,parsep=0pt,leftmargin=1.45em}
\setlist[enumerate]{topsep=2pt,itemsep=2pt,parsep=0pt,leftmargin=1.65em}
\setcounter{secnumdepth}{0}
\renewcommand{\arraystretch}{1.12}

% Zahlavi: vlevo rocnik | globalni tema, vpravo cislo: nazev hodiny
% Zapati zustava zcela prazdne.
\pagestyle{fancy}
\fancyhf{}
\renewcommand{\headrulewidth}{0.35pt}
\renewcommand{\footrulewidth}{0pt}
\setlength{\headheight}{14pt}
\newcommand{\nastavPripravu}[4]{%
  \fancyhead[L]{\small #1\enspace\textbar\enspace #2}%
  \fancyhead[R]{\small #3: #4}%
}

% Minimalisticke cernobile boxy. Kazda definice zvlast.
\tcbset{commonlesson/.style={enhanced,breakable=false,colback=white,colframe=black,
  colbacktitle=white,coltitle=black,boxrule=0.45pt,
  arc=0.8mm,outer arc=0.8mm,left=3mm,right=3mm,top=1.6mm,bottom=1.6mm,
  toptitle=1mm,bottomtitle=0.7mm,fonttitle=\bfseries,
  attach boxed title to top left={xshift=2.5mm,yshift=-1.15mm},
  boxed title style={colback=white,colframe=white,boxrule=0pt,left=1mm,right=1mm}}}
\newtcolorbox{zadaniBox}[1][]{commonlesson,title={Zadání},#1}
\newtcolorbox{prikladBox}[1][]{commonlesson,title={Příklad},#1}
\newtcolorbox{definiceBox}[1][]{commonlesson,title={Definice},#1}

% Spolecny box pro zadani a jeho reseni, bez mezery mezi castmi.
% Pouziti:
%   \begin{prikladReseniBox}
%     Zadani ulohy.
%     \tcblower
%     \begin{align*} ... \end{align*}
%   \end{prikladReseniBox}
% Spodni cast ma jemny sedy podklad a delici linku; lze v ni sazet
% vice radku matematiky. Obe casti maji jen jeden spolecny ramecek.
\tcbset{solutionlesson/.style={
  commonlesson,
  enhanced jigsaw,
  breakable,
  colbacklower=black!3,
  segmentation style={solid,draw=black!25,line width=0.35pt},
  fontlower=\small\color{notegray},
  before lower={\noindent\textbf{Řešení:}\enspace},
  leftlower=3mm,rightlower=3mm
}}
\newtcolorbox{prikladReseniBox}[1][]{solutionlesson,title={Příklad},#1}
\newtcolorbox{zadaniReseniBox}[1][]{solutionlesson,title={Zadání},#1}

% Pomucky pro autora pripravy: sedive poznamky; \reseni pro starsi pripravy
\newcommand{\poznamka}[1]{\par\begingroup\small\color{notegray}\itshape
  \textnormal{\textbf{Poznámka pro učitele.}}\enspace #1\par\endgroup}
\newcommand{\reseni}[1]{\par\begingroup\small\color{notegray}
  \textbf{Očekávané řešení:}\enspace #1\par\endgroup}
\newcommand{\blok}[1]{\needspace{5\baselineskip}\subsection*{#1}}
\newcommand{\vyrok}[1]{\uv{#1}}


\nastavPripravu{1. ročník}{Mocniny}{25}{Mocnina s přirozeným exponentem}

\begin{document}

\blok{Výklad a souvislosti}

Zápis $4+4+4$ zkracujeme násobením, zápis $4\cdot4\cdot4$ mocninou. Exponent tedy neříká, čím základ násobíme, ale kolikrát se opakuje jako činitel. Dnes používáme jen kladné celé exponenty; nula a záporné exponenty přijdou ve 28. hodině.

\begin{definiceBox}[title={Definice},breakable=false]
Pro $a\in\mathbb R$ a $n\in\mathbb N=\{1,2,\ldots\}$ je \textbf{mocnina s přirozeným exponentem} $a^n=\underbrace{a\cdot a\cdots a}_{n\text{ činitelů}}$.
\end{definiceBox}


\begin{definiceBox}[title={Definice},breakable=false]
\textbf{Základ mocniny} je číslo $a$ v zápisu $a^n$, které opakovaně násobíme.
\end{definiceBox}


\begin{definiceBox}[title={Definice},breakable=false]
\textbf{Exponent} $n$ v zápisu $a^n$ udává při $n\in\mathbb N$ počet stejných činitelů $a$.
\end{definiceBox}


Umocňování má přednost před násobením znaménkem minus. U záporného základu rozhoduje sudost exponentu: dvojice záporných činitelů mají kladný součin. Proto $(-a)^{2k}=a^{2k}$ pro $k\in\mathbb N$ a $(-a)^{2k+1}=-a^{2k+1}$ pro celá $k\ge0$. Dále $a^1=a$, $1^n=1$ a $0^n=0$.

\blok{Řešené příklady}

\begin{prikladReseniBox}[title={Příklad 1},breakable=false]
Porovnej $(-2)^4$, $-2^4$ a $(-2)^3$.
\tcblower
$(-2)^4=16$, zatímco $-2^4=-(2^4)=-16$. Dále $(-2)^3=-8$. V prvním zápisu se mocní celý záporný základ.
\end{prikladReseniBox}


\begin{prikladReseniBox}[title={Příklad 2},breakable=false]
Vypočítej $(-3)^3$ a $\left(\frac12\right)^5$.
\tcblower
$(-3)^3=(-3)(-3)(-3)=-27$; $\left(\frac12\right)^5=\frac1{32}$. Součin pěti polovin není $\frac52$.
\end{prikladReseniBox}


\blok{Samostatné procvičení}

\begin{zadaniReseniBox}[title={Procvičení},breakable=false]
\begin{enumerate}[label=\alph*)]
\item Vypočítej $300^1$, $0^7$, $2^5$ a $1^{12}$.
\item Vypočítej $(-1)^{35}$ a rozhodni, zda $(-5)^{15}>0$.
\item Vypočítej $-3^2+(-3)^2$.
\end{enumerate}
\tcblower
\begin{enumerate}[label=\alph*)]
\item Postupně $300$, $0$, $32$, $1$; použijeme počet činitelů z definice.
\item $(-1)^{35}=-1$. Tvrzení je nepravdivé: záporný základ na lichý exponent dává záporné číslo.
\item $-9+9=0$. Závorky mění základ mocniny.
\end{enumerate}
\end{zadaniReseniBox}


\poznamka{Nech nejprve žáky rozepsat několik součinů. Sudost zdůvodni párováním znamének, nikoli jen tabulkou pokusů. Pojmy základ a exponent vysvětli odděleně. Zatím nepoužívej vzorce pro mocniny součtu.}


\blok{Navazující procvičování ve třídě}

\begin{zadaniReseniBox}[title={Základ: závorky a pořadí operací},breakable=false]
\begin{enumerate}[label=\alph*)]
\item Vypočítej $(-4)^2$, $-4^2$, $-(-4)^2$, $(-1)^{2026}$. U každého výrazu nejprve pojmenuj základ mocniny.
\item Vypočítej $(-2)^3+3^2$, $2\cdot(-3)^2-4^2$, $(-2+5)^3$ a $-(-2)^3$.
\item Vypočítej $(-\tfrac12)^4$, $(-\tfrac12)^3$ a $(\tfrac32)^2$.
\end{enumerate}
\tcblower
\begin{enumerate}[label=\alph*)]
\item $16$, $-16$, $-16$, $1$. V prostředních dvou výrazech je mínus před mocninou; poslední mínus v $-(-4)^2$ tedy nemizí.
\item $1$, $2$, $27$, $8$. Nejprve vnitřní závorky, pak mocniny, násobení a sčítání.
\item $\frac1{16}$, $-\frac18$, $\frac94$; jde o násobení čtyř, tří, resp. dvou stejných činitelů.
\end{enumerate}
\end{zadaniReseniBox}

\begin{zadaniReseniBox}[title={Základ a navazující úvaha: porovnávání},breakable=false]
\begin{enumerate}[label=\alph*)]
\item Doplň $<,=,>$: $(-3)^4\ \ldots\ 3^4$; $-3^4\ \ldots\ (-3)^4$; $(-2)^7\ \ldots\ (-2)^8$.
\item Porovnej $(\tfrac12)^5$ a $(\tfrac12)^4$. Vysvětli bez desetinného vyčíslení.
\item Rozhodni a zdůvodni: $(-a)^6=a^6$ a $(-a)^5=a^5$ pro všechna reálná $a$.
\end{enumerate}
\tcblower
\begin{enumerate}[label=\alph*)]
\item Po řadě $=$, $<$, $<$. U posledního porovnání stačí rozdílná znaménka.
\item $\frac1{32}<\frac1{16}$. Další násobení číslem $\frac12$ kladnou hodnotu zmenší.
\item První rovnost platí, protože se šest znamének minus spáruje. Druhá obecně neplatí, například pro $a=1$ je $-1\ne1$; správně $(-a)^5=-a^5$.
\end{enumerate}
\end{zadaniReseniBox}

\begin{zadaniReseniBox}[title={Rezerva: propojené výpočty},breakable=false]
\begin{enumerate}[label=\alph*)]
\item Vypočítej $2^3-(-2)^3+(-3)^2-3^2$ a $\bigl((-2)^2-3\bigr)^5$.
\item Porovnej $(2-5)^4$ a $(5-2)^4$, pak $(2-5)^3$ a $(5-2)^3$.
\end{enumerate}
\tcblower
\begin{enumerate}[label=\alph*)]
\item $8-(-8)+9-9=16$; druhý výraz je $(4-3)^5=1$.
\item Sudé mocniny jsou obě $81$, liché jsou $-27$ a $27$, tedy první je menší. Opačná čísla mají shodné sudé mocniny.
\end{enumerate}
\end{zadaniReseniBox}


\poznamka{Původní výklad a příklady jsou úvod. Ve třídě navazují obě základní skupiny; u každé požaduj nejprve samostatný zápis, potom vysvětlení znaménka či závorky. Pro rychlé slouží propojené výpočty, domácí úlohy jsou oddělené. Používáme jen definici a pořadí operací, nikoli pravidla příštích hodin.}

\blok{Domácí procvičení}

\begin{zadaniReseniBox}[title={Domácí úlohy},breakable=false]
\begin{enumerate}[label=\alph*)]
\item Vypočítej $(-5)^2$, $-5^2$, $-(-5)^3$.
\item Vypočítej $(-\tfrac23)^3$ a $(\tfrac14)^2$.
\item Urči $0^{15}$, $1^{25}$ a $(-1)^{25}$.
\item Vypočítej $3^2-2\cdot(-2)^3$.
\item Porovnej $(-4)^3$ a $-4^3$; $(\tfrac13)^2$ a $(\tfrac13)^3$.
\item Najdi chybu v zápisu $(-2+3)^2=(-2)^2+3^2$ a spočítej obě strany.
\end{enumerate}
\tcblower
\begin{enumerate}[label=\alph*)]
\item $25$, $-25$, $125$.
\item $-\frac8{27}$ a $\frac1{16}$.
\item $0$, $1$, $-1$.
\item $9-2\cdot(-8)=25$.
\item První dvojice se rovná ($-64$); ve druhé $\frac19>\frac1{27}$.
\item Levá strana je $1$, pravá $13$. Mocnina součtu se takto nerozděluje.
\end{enumerate}
\end{zadaniReseniBox}

\end{document}
